% !TeX root = ../thuthesis-example.tex

\chapter{总结与展望}

\section{全文工作总结}

本论文围绕量子几何在拓扑物态与非线性输运中的多层面效应，沿"几何对象—外场耦合—可观测信号"的主线展开三部分研究。量子几何张量作为布洛赫态在参数空间中变化快慢的内禀度量，其实部——量子度规——与虚部——贝里曲率——分别编码了布洛赫态之间的距离与旋转信息。本论文的核心目标，是阐明这两个互补的几何量如何在从静态拓扑分类到弱场响应、再到强场非平衡动力学的不同物理层面产生可观测效应。

\subsection{广义嵌套Wilson环与高阶拓扑诊断}

第~\ref{chap:hot-generic-nwl}~章将嵌套Wilson环理论从镜面对称体系推广到包含滑移反射在内的一般晶体对称情形。传统嵌套Wilson环依赖于镜面反射对Wannier能带的量子化约束，但大量空间群——特别是包含非共形对称性（如滑移反射）的体系——并不满足这一前提。为此，本章引入广义Wannier扇区的划分方案，允许Wannier能带在对称性约束下分裂为$N_s \ge 2$个扇区，并定义各扇区上的广义Wannier扇区极化（gNBP）作为高阶拓扑不变量。

在此框架下，本章对80个二维层群中相关的对称性约束进行系统分析，识别出9个支持传统嵌套Berry相$\theta^{(2)}$和15个支持广义嵌套Berry相$\theta^{(4)}$的层群。在材料层面，本章以NaCuSe和KMgBi两个狄拉克半金属为原型体系，利用第一性原理紧束缚模型验证了广义嵌套Berry相与高阶费米弧之间的对应关系。特别值得注意的是，在KMgBi中，传统的$\theta^{(2)}$未能捕获高阶拓扑相变，而$\theta^{(4)}$成功给出了与棱态谱一致的诊断结果——这一对比有力地论证了推广嵌套Wilson环的必要性。

本章工作的方法论意义在于：它将高阶拓扑的诊断工具从一小类具有镜面对称性的体系扩展到更广泛的晶体结构中，使得"对称性约束$\Rightarrow$Wannier能带量子化$\Rightarrow$高阶不变量$\Rightarrow$边界谱预言"的推导链条具有更强的普适性。

\subsection{弱场非线性响应中的量子几何对照}

第~\ref{chap:weak-field}~章转入弱场低频条件下的二阶横向输运，核心问题是：在给定对称性破缺方式下，量子几何张量的不同分量如何分别映射为可测量的非线性霍尔信号。

本章沿两条互补的机制链展开。第一条围绕量子度规，以单层反铁磁体MnS为平台，揭示了$\mathcal{PT}$对称体系中由贝里联络极化率（BCP）偶极子驱动的本征非线性霍尔效应（INH）\cite{PhysRevLett.131.056401}。尽管$\mathcal{PT}$对称严格禁戒贝里曲率与BCD，量子度规仍可通过BCP偶极子提供非零的本征响应，其电导率随Néel矢量方向$\theta$的变化为反铁磁序参量的全电学探测提供了方案。第二条围绕贝里曲率偶极子（BCD），以ZrTe$_5$薄膜为平台，揭示了层间滑动诱导的自发反演对称性破缺及其引起的非线性反常霍尔效应（NAHE）与动力学磁电效应（KME）\cite{nsr_nwad103}。$k\cdot p$模型的解析推导明确建立了滑动位移$\lambda$与BCD/KME之间的线性对应关系，揭示了倾斜项在生成非互易响应中的主导作用；多层分析表明这一机制并非三层薄膜的特例，而是可推广至不同厚度薄膜与自然解理表面的普遍现象。

两条机制链的物理对比深化了对量子几何框架的理解：BCD机制作为费米面性质与$\tau$成正比，INH机制作为费米海性质独立于$\tau$。二者可通过频率依赖关系在实验中加以区分，也可能在同一材料平台中竞争或协同。这一对照揭示了量子几何张量的两个分量——贝里曲率与量子度规——在弱场非线性输运中扮演的互补角色：前者主导反演对称性破缺体系中的BCD信号，后者在时间反演对称但破缺其他对称的体系中提供独立的响应通道。

\subsection{强场量子度量诱导振荡}

第~\ref{chap:strong-field-qgeom}~章将研究推进到强电场驱动下的非平衡量子动力学。在强场条件下，传统的微扰展开不再适用，本章提出以大能隙近似为核心的密度矩阵递推框架，将演化系统地分解为对角占据数（零阶）与带间相干（高阶）的层级结构，统一给出漂移电流与振荡电流的表达式。

基于该框架，本章提出并刻画了量子度量诱导振荡（QMO）这一新的几何主导输运机制。QMO的物理根源在于：量子度规在布里渊区中的空间涨落使得强场驱动下的带间相干贡献呈现周期性振荡，其频率由能隙结构决定，但振幅与量子度规的动量依赖直接相关。QMO具有两个关键特征：其一，它可在严格平带中存在——这与依赖色散的布洛赫振荡（BO）形成本质区别；其二，在一定对称性条件下，QMO振荡分量的幅度不随电场增强而衰减，而BO振幅在强场下被抑制。这些区分特征为实验识别提供了判据。在数值层面，本章通过一维GaAs/AlGaAs超晶格模型和二维蜂窝晶格模型验证了上述解析预言，展示了QMO的特征频谱和漂移标度；在严格平带模型中，QMO通过呼吸模信号直观呈现，其图样不随电场强度变化，与BO形成鲜明对比。

在弱场极限下，递推框架自洽地还原为第~\ref{chap:weak-field}~章讨论的非线性响应结果：BCD机制对应$\tau$线性的时间反演偶响应，INH机制对应$\tau$无关的时间反演奇响应。这一还原关系的物理意义在于：弱场非线性霍尔效应与强场QMO并非两套独立的理论，而是同一组递推方程在不同电场尺度下的不同极限。换言之，量子几何在输运中的表现形式——从弱场的响应系数到强场的动力学信号——构成一个连续谱，本章的递推框架为描述这一连续过渡提供了统一工具。

\subsection{核心贡献与统一图像}

综合第~\ref{chap:hot-generic-nwl}--\ref{chap:strong-field-qgeom}~章的结果，本论文的贡献可从以下层面理解。

\begin{enumerate}
  \item \textbf{方法层面。}建立广义嵌套Wilson环的可计算框架并完成层群分类；构造弱场二阶响应的统一表达并区分度规/曲率贡献的对称性选择规则；提出强场大能隙递推方法并导出涵盖漂移与振荡的解析输运表达式。
  \item \textbf{机制层面。}形成"对称性—量子几何—可观测响应"的贯穿链条。静态层面，量子几何通过Wilson环编码高阶拓扑不变量，决定边界谱结构；弱场层面，量子度规与贝里曲率通过不同的响应通道映射为非线性霍尔信号，可作为内部序参量的电学读出；强场层面，量子度规的动量空间涨落驱动新型振荡机制（QMO），将几何信息印刻在动力学信号中。三者由量子几何张量这一共同的数学对象联结。
  \item \textbf{面向材料与实验。}在若干具体平台给出可检验的预言与判据：狄拉克半金属中的高阶费米弧、反铁磁薄膜中基于INH的Néel矢量读出、ZrTe$_5$薄膜中滑动诱导的NAHE信号、以及超晶格与莫尔体系中QMO的特征频谱与场强标度。
\end{enumerate}

上述三个层面中，"对称性—量子几何—可观测响应"链条构成本论文的核心洞察。从拓扑量子化学（第~\ref{chap:hot-generic-nwl}~章）到弱场非线性输运（第~\ref{chap:weak-field}~章）再到强场非平衡动力学（第~\ref{chap:strong-field-qgeom}~章），晶体对称性始终是决定量子几何能否进入可观测通道的选择规则来源；而量子几何张量的实部与虚部则分别在不同的物理条件下担当主导角色。这一图像的统一性不仅体现在物理概念层面——三章工作共享量子几何张量这一数学对象——还体现在技术层面：第~\ref{chap:strong-field-qgeom}~章的递推框架在弱场极限下精确还原为第~\ref{chap:weak-field}~章的响应表达式，从而将两章的结果纳入同一组方程的不同极限。

需要指出的是，本文以解析理论与模型推导为主，对无序散射、电子关联、有限温度以及弛豫过程的系统纳入仍有提升空间。部分预测信号的观测窗口依赖于实验条件（电场强度、散射时间、样品均匀性等）的精细控制。对这些因素的定量处理，是将几何机制从理论框架推进到可直接比对实验的关键步骤。

\section{研究展望}

以上三个方向的工作虽各有侧重，但它们提示了若干具有内在联系的后续研究方向。以下的展望基于论文各章中已有的讨论与分析，从各章的自然延伸与跨章节的贯通方向两个层面展开。

\subsection{高阶拓扑：从不变量到可观测边界谱}

第~\ref{chap:hot-generic-nwl}~章建立的不变量框架为高阶拓扑的诊断提供了工具，但从不变量到实验验证仍需若干关键步骤。首先，在候选材料（如NaCuSe、KMgBi及其同构体系）中建立"结构与对称性控制—边界谱观测—不变量验证"的闭环是当务之急：ARPES和STM等谱学手段可用于锁定高阶费米弧的能量与空间分布，而表面/界面对称性的可控调节（如应变、化学修饰）则允许检验保护机理。其次，广义嵌套Wilson环形式有望推广到更广泛的对称性设置中，为脆弱拓扑\cite{PhysRevLett.123.186401,PhysRevB.100.115160,PhysRevB.100.205126}等近年涌现的概念提供统一的几何诊断。面向应用，高阶拓扑边界态（角态、铰链态）的局域化特性使其成为潜在的读出或耦合单元，但将拓扑判据转化为器件设计规则，需要在有效模型与第一性原理参数之间建立可追踪的映射。

\subsection{弱场非线性响应：分量分离与材料筛选}

第~\ref{chap:weak-field}~章的对照式刻画揭示了BCD与INH机制的对称性选择规则和弛豫时间依赖的差异，这为后续研究提供了两个互补的推进方向。

在材料侧，可以对称性破缺方式为索引，结合高通量计算筛选具有大BCD或强量子度规不均匀性的候选体系，形成可复用的筛选流程。ZrTe$_5$薄膜中揭示的"低能结构自由度驱动量子几何响应"这一机制，提示在更多范德瓦尔斯薄膜与表面体系中系统识别自发对称性破缺的可能性——滑动铁电性、层间扭转等低能模式可能提供新的$\mathcal{I}$破缺来源，其与非线性输运信号的定量对应关系值得建立。在实验侧，利用角向依赖、频率依赖与对称性选择规则区分度规项与曲率项，并系统抑制热效应与寄生非线性信号，是实现分量分离的关键。将二阶响应推广到交流与太赫兹频段，有望把量子几何信号从材料表征推进到功能器件的设计。

此外，在关联、磁性或莫尔超晶格体系中，量子度规与贝里曲率对不同非线性响应通道的竞争与协同尚待系统研究。这一方向不仅涉及对弛豫时间近似的超越——需要系统处理侧跳\cite{Du2021NonlinearHE}、斜散射\cite{Du2021NonlinearHE}与反常斜散射\cite{PhysRevLett.131.076601}等散射机制对非线性输运的修正——也与第~\ref{chap:strong-field-qgeom}~章讨论的强场理论存在方法论上的联系：后者附录中给出的一般弛散算符下的递推公式，为在统一框架中处理散射效应提供了出发点。

\subsection{强场动力学：QMO的观测窗口与理论完善}

第~\ref{chap:strong-field-qgeom}~章的QMO预言为强场区间的量子几何识别提供了新判据，但将其推进到可重复观测需要在实验与理论两个层面同时努力。

在实验侧，一维外延超晶格长期是研究BO与强场输运的成熟平台\cite{LEO1992943,PhysRevLett.64.3167}，本章结果提示其中观测到的退局域化电流可能包含量子度规驱动的带间相干贡献，值得通过对电场标度关系的系统测量加以重新审视。二维莫尔超晶格具有可调平带与丰富量子几何信息\cite{lau2022reproducibility,Bloch_in_Moire,Quantum_Geometric_Oscillations}，是比较$J_{\text{Bloch}}$、$J_{\Omega}$与$J_g$贡献的理想平台。光学晶格与冷原子体系在时间分辨密度测量方面的优势\cite{PhysRevLett.87.140402,PhysRevLett.76.4508}，则使其成为通过呼吸模探测量子度规涨落的候选实验方案——QMO振幅不随电场衰减的特征为区分信号来源提供了关键判据。

在理论侧，第~\ref{chap:strong-field-qgeom}~章已指出若干需要系统解决的开放问题。第一，大能隙近似的适用边界：当电场接近齐纳击穿条件$eEa \sim \Delta^2/W$时，带间隧穿的非微扰重分布使递推展开失效；在平带系统中，带宽$W$趋零使齐纳击穿条件发生质变，量子几何诱导的击穿机制可能呈现新的物理特征，需要与WKB方法或Dykhne-Davis-Pechukas公式相衔接\cite{Kitamura2019NonreciprocalLT,L-Z_Green_RPB_2020}。第二，超越弛豫时间近似：更一般的弛散机制——包括杂质散射导致的侧跳与斜散射\cite{Du2021NonlinearHE}、反常斜散射\cite{PhysRevLett.131.076601}——在弱场非线性输运中已被证明贡献显著，但它们在强场条件下的行为尚未被系统研究。第三，电子关联效应：在莫尔超晶格等平带体系中，相互作用能标$V$与动能可比\cite{cao2018unconventional,cao2018correlated,cai2023signatures}，当$V \sim eEa$时电场不仅驱动单粒子动力学还可能与多体关联态竞争或协同，将单粒子量子几何框架推广到包含相互作用的多体理论是连接QMO与强关联实验的必要步骤。第四，有限温度下退相干对BO与QMO信号信噪比的定量影响——特别是声子散射与电子-电子散射引入的$\tau$温度依赖——将为实验方案的优化提供指导。

\subsection{贯通方向：从静态拓扑到非平衡几何的统一视角}

本论文三部分工作分别处于平衡拓扑分类、弱场非线性响应和强场非平衡动力学三个层面，但它们由量子几何张量这一核心数学对象串联。更具前瞻性的后续方向，是将这三个层面置于一套连续的描述框架中。

一个具体的推进路径是建立弱场到强场的连续标度描述。第~\ref{chap:strong-field-qgeom}~章已展示递推框架在弱场极限下自洽还原为第~\ref{chap:weak-field}~章的非线性响应结果。在此基础上，构建中等场强区间（$eEa$与$\hbar/\tau$可比但仍小于$\Delta$）的标度律与拟合判据，将为实验上在同一样品中连续追踪从弱场响应系数到强场动力学信号的过渡提供理论支撑。

另一个值得探索的方向，是将高阶拓扑边界态的局域化特性与量子几何输运读出机制结合。第~\ref{chap:hot-generic-nwl}~章讨论的角态和铰链态具有受对称性保护的空间局域化，而第~\ref{chap:weak-field}~章展示的非线性霍尔效应对局域量子几何分布敏感。二者的结合提示了一种可能的"拓扑保护+几何信号放大"器件概念——利用边界态的拓扑保护确保信号的鲁棒性，同时利用其增强的量子几何（如大曲率或强度规不均匀性）放大非线性响应。这一设想的可行性需要在具体材料平台上通过第一性原理与输运计算加以评估。

最后，发展连接第一性原理电子结构、有效模型与输运计算的多尺度工作流程，是支撑从材料筛选到器件性能预测的闭环设计的基础设施层面的需求。本文使用的HopTB\cite{HopTB}软件包为此提供了初步的技术基础——它支持拓扑、输运和响应性质的计算，并具有与Wannier90\cite{Wannier90v2}和OpenMX\cite{Openmx2003,Openmx2004}的接口——但在自动化程度、多体效应的纳入以及与实验数据的对接方面仍有较大的发展空间。
